\section{Justificación}

    El análisis de WebSIS resulta pertinente porque impacta directamente en la experiencia diaria de miles de estudiantes de la UMSS. Al ser el único canal institucional habilitado para la gestión académica, sus deficiencias afectan de manera transversal a toda la población estudiantil, con mayor intensidad durante los periodos de inscripciones y publicación de calificaciones, cuando los accesos se concentran y las consecuencias de los errores son más graves.

    Desde la perspectiva de la Interacción Humano-Computadora, el proyecto aplica directamente los principios establecidos por \textcite{nielsen1994}, cuyas diez heurísticas de usabilidad —entre ellas la visibilidad del estado del sistema, la correspondencia entre el sistema y el mundo real, y el control y libertad del
    usuario— son precisamente las dimensiones donde WebSIS presenta sus mayores deficiencias. Del mismo modo, los principios de diseño centrado en el usuario formulados por \textcite{norman2002} sostienen que un sistema interactivo no puede considerarse funcional si el usuario no logra operarlo de manera eficiente, autónoma y sin errores, con independencia de su corrección técnica interna. WebSIS cumple sus funciones desde un punto de vista técnico, pero falla sistemáticamente en estas dimensiones de
    interacción.

    La urgencia del estudio se sostiene en el carácter obligatorio de los procesos que la plataforma gestiona. Los estudiantes ingresantes deben adaptarse simultáneamente a la vida universitaria y a una interfaz poco intuitiva, asumiendo una doble carga cognitiva en el momento de mayor vulnerabilidad dentro de su trayectoria académica. Los estudiantes regulares, por su parte, han construido rutinas de uso que compensan las deficiencias del sistema, pero esa adaptación individual no elimina el problema de diseño subyacente: evidencia que el sistema delega en el usuario una carga que debería asumir la interfaz.

    Existe, en consecuencia, una necesidad real y documentada de analizar WebSIS desde los principios de la IHC y de proponer mejoras concretas que transformen su experiencia de uso, colocando al estudiante en el centro del proceso de diseño de la interacción.